\begin{helper}[TwoBiteRedMarginal]
\label{helper:TwoBiteRedMarginal}
Let \(n \le m^2\). For any event \(E\) depending only on the red graph, the probability of \(E\) under \(\noderef{TwoBiteEventProbability}(n, m, p)\) is equal to its marginal probability under \(\noderef{GnpGraphWeight}(p, \cdot)\):
\[
\Pr(E(\noderef{TwoBiteRedGraph}(\omega))) = \sum_{G : E(G)} \noderef{GnpGraphWeight}(p, G).
\]
\end{helper}

\begin{proof}
By \(\noderef{TwoBiteSample}\), a sample is a triple
\[
\omega=(G_R,G_B,\pi),
\]
where \(G_R\) is the red base graph, \(G_B\) is the blue base graph, and
\(\pi:\operatorname{Fin}(n)\hookrightarrow
\operatorname{Fin}(m)\times\operatorname{Fin}(m)\) is the injection.  The
definition \(\noderef{TwoBiteRedGraph}\) extracts the first coordinate
\(G_R\).

Expand the event probability using \(\noderef{TwoBiteEventProbability}\) and
\(\noderef{TwoBiteSampleWeight}\):
\[
\Pr(E(\noderef{TwoBiteRedGraph}(\omega)))
=
\sum_{G_R,G_B,\pi}
\mathbf 1_{E(G_R)}\,
\noderef{GnpGraphWeight}(p,G_R)\,
\noderef{GnpGraphWeight}(p,G_B)\,
\noderef{UniformInjectionWeight}(n,m).
\tag{1}
\]
Fix the projected red graph \(G_R\).  In the inner sum over the irrelevant
coordinates \(G_B\) and \(\pi\), the factor
\(\mathbf 1_{E(G_R)}\noderef{GnpGraphWeight}(p,G_R)\) is constant.  Since
\(n\le m^2\), the injection type is nonempty, so
\(\noderef{UniformInjectionWeight}(n,m)\) is the reciprocal of the number of
injections.  Therefore
\[
\sum_\pi \noderef{UniformInjectionWeight}(n,m)=1.
\tag{2}
\]
The blue graph coordinate has total mass
\[
\sum_{G_B}\noderef{GnpGraphWeight}(p,G_B)=1
\tag{3}
\]
by \(\noderef{GnpGraphWeightSumOne}\).  Substituting (2) and (3) into (1)
leaves
\[
\sum_{G_R}
\mathbf 1_{E(G_R)}\,\noderef{GnpGraphWeight}(p,G_R),
\]
which is exactly
\[
\sum_{G_R:\,E(G_R)} \noderef{GnpGraphWeight}(p,G_R).
\]
This is the asserted red marginal identity.
\end{proof}
